\chapter{Métodos intrínsecamente semi-supervisados}

En este capítulo se describen los métodos intrínsecamente semi-supervisados, es decir, aquellos que desde su concepción están diseñados para entrenar de manera conjunta con datos etiquetados y no etiquetados. Su rasgo distintivo es que los ejemplos no etiquetados no aparecen como una etapa auxiliar o previa, sino que se incorporan de forma explícita en la formulación del método: bien a través de un modelo probabilístico (por ejemplo, en la verosimilitud de $p(\mathbf{x},y)$), o bien como términos adicionales en la función objetivo (por ejemplo, regularización, consistencia o suavidad).

Siguiendo la última división de los métodos inductivos presentada en la taxonomía del Capítulo 3, organizamos los métodos intrínsecos en tres grandes grupos: métodos generativos, métodos discriminativos y métodos basados en variedades. En cada caso se introduce la intuición del enfoque y se presentan formulaciones representativas: mezclas de modelos optimizadas con EM, clasificadores de máximo margen como las S3VM y métodos basados en perturbaciones como las \textit{ladder networks}, y, finalmente, técnicas de regularización de la variedad mediante grafos.

\section{Generativos}

Los métodos generativos constituyen uno de los enfoques clásicos del aprendizaje semi-supervisado \cite{zhu2005semi}. Su idea central es modelar explícitamente el proceso que genera los datos, esto es, la distribución conjunta $p(\mathbf{x},y)$, y, una vez estimada, usarla de diferentes formas para clasificar \cite{van2020survey}. Este enfoque conecta directamente con la discusión del Capítulo 2 sobre cómo la información de $p(\mathbf{x})$ puede ayudar a inferir $p(y\mid\mathbf{x})$ (Sección \ref{sec:suposiciones}).

Dentro de esta familia, el referente clásico son los modelos de mezcla optimizados mediante máxima verosimilitud, normalmente con el algoritmo EM cuando hay variables latentes \cite{dempster1977maximum}.

\subsection{Mezcla de modelos}

Los modelos de mezcla asumen que los datos se generan a partir de una combinación de componentes probabilísticas, típicamente una por clase, con parámetros desconocidos \cite{van2020survey}. Bajo este supuesto, los datos etiquetados y no etiquetados pueden contribuir conjuntamente a estimar el modelo.

La intuición en semi-supervisado es que los datos no etiquetados ayudan a identificar la estructura de la mezcla (componentes, pesos y forma), mientras que los pocos datos etiquetados permiten asociar componentes a clases \cite{zhu2005semi,zhu2009introduction}.

\subsubsection{Mezcla de modelos supervisada}

Dada una instancia $\mathbf{x}$, la predicción se obtiene tomando la clase más probable a posteriori:

\begin{equation}
    \hat{y} = \operatorname*{arg\,max}_{y} p(y|\mathbf{x})
\end{equation}

Para realizar dicho cálculo, se emplea la regla de Bayes:

\begin{equation}
    p(y|\mathbf{x}) = \frac{p(\mathbf{x}|y)p(y)}{\sum_{y'}p(\mathbf{x}|y')p(y')}
\end{equation}

por lo que el problema se reduce a estimar $p(\mathbf{x}|y)$ y $p(y)$. En el contexto de una mezcla de componentes probabilísticos, esto equivale a estimar los parámetros que definen dichas componentes.  

Si denotamos por $\theta$ al vector con los parámetros que definen los modelos y por $D_l=\{(\mathbf{x}_i,y_i)\}_{i=1}^{l}$ al conjunto etiquetado. Entonces el estimador de máxima verosimilitud de los parámetros es

\begin{equation}
    \hat{\theta} = \operatorname*{arg\,max}_{\theta} p(D_l|\theta) = \operatorname*{arg\,max}_{\theta} \log p(D_l|\theta).
\end{equation}

Es decir, buscamos los parámetros $\theta$ que maximizan la verosimilitud de los datos de entrenamiento. Además, como estamos tratando con datos independientes e idénticamente distribuidos, la verosimilitud se factoriza como el producto de las probabilidades individuales, lo que nos lleva a la siguiente expresión para la log-verosimilitud:

\begin{equation}
    \log p(D_l|\theta) = \log \prod_{i=1}^{l} p(\mathbf{x}_i, y_i|\theta)
     = \sum_{i=1}^{l} \log \bigl(p(y_i|\theta)p(\mathbf{x}_i|y_i,\theta)\bigr).
\end{equation}

Para familias concretas (por ejemplo, mezclas gaussianas con etiquetas observadas), este problema suele admitir estimadores cerrados o de optimización relativamente directa \cite{zhu2009introduction}.

\subsubsection{Mezcla de modelos semi-supervisada}

En el caso semi-supervisado se dispone, además de $D_l$, de un conjunto no etiquetado $D_u=\{\mathbf{x}_j\}_{j=l+1}^{l+u}$. El conjunto completo es $D=D_l\cup D_u$, y la log-verosimilitud pasa a ser:

\begin{equation}
    \log p(D|\theta) = \sum_{i=1}^{l} \log p(\mathbf{x}_i,y_i|\theta) + \sum_{j=l+1}^{l+u} \log p(\mathbf{x}_j|\theta).
\end{equation}

Donde en el término de los datos no etiquetados se incorpora la distribución marginal
\begin{equation}
    p(\mathbf{x}_j|\theta)=\sum_{y\in\mathcal{Y}} p(\mathbf{x}_j,y|\theta)=\sum_{y\in\mathcal{Y}} p(y|\theta)p(\mathbf{x}_j|y,\theta),
\end{equation}
que refleja que la etiqueta de $\mathbf{x}_j$ es desconocida \cite{zhu2009introduction}.

En general, esta optimización deja de ser cóncava y no suele resolverse de forma analítica. Por ello, el método estándar es EM, un procedimiento iterativo de máxima verosimilitud con variables latentes \cite{dempster1977maximum}. De forma intuitiva, EM repite el siguiente ciclo:
\begin{itemize}
    \item \textbf{Inicialización}: partir de unos parámetros iniciales $\theta^{(0)}$ (por ejemplo, obtenidos con los datos etiquetados).
    \item \textbf{E-step}: con los parámetros actuales, estimar para cada dato no etiquetado la probabilidad de pertenecer a cada clase, es decir, una etiqueta \textit{blanda}.
    \item \textbf{M-step}: de manera análoga al caso supervisado, reestimar los parámetros del modelo, pero usando ahora las etiquetas probabilísticas obtenidas en la E-step como pesos en la máxima verosimilitud.
    \item \textbf{Iteración}: repetir E-step y M-step hasta convergencia.
\end{itemize}

Esta dinámica es conceptualmente parecida al pseudoetiquetado de los métodos envolventes (Capítulo \ref{cap:envolventes}) \cite{zhu2009introduction,zhou2010semi}, pero aquí las asignaciones para $D_u$ son probabilísticas, mientras que en los métodos envolventes se termina obteniendo una etiqueta final para cada ejemplo.

La mezcla de modelos es un enfoque que puede ser muy eficaz cuando el modelo generativo está bien especificado. Sin embargo, su rendimiento depende críticamente de hipótesis como la identificabilidad de la mezcla y la corrección del modelo; si estas fallan, los datos no etiquetados pueden degradar el rendimiento \cite{zhu2005semi,van2020survey}.


\section{Discriminativos}

Esta familia de métodos se distingue de la anterior en que no trata de modelar el proceso generativo ni la densidad conjunta $p(\mathbf{x},y)$, sino que se centra directamente en aprender un modelo discriminativo $p(y\mid\mathbf{x})$ o una función de decisión $f:\mathcal{X}\to\mathcal{Y}$ \cite{ng2001discriminative}. En este caso, la información de los datos no etiquetados se incorpora con el objetivo de mejorar la generalización del modelo, normalmente a través de suposiciones adicionales sobre la distribución subyacente de los datos (por ejemplo, la separación en baja densidad o la suposición de suavidad) \cite{van2020survey}.
\subsection{Máximo margen}
\label{subsec:maximo_margen}

Los métodos de máximo margen buscan aprender una frontera de decisión que quede lo más alejada posible de los ejemplos de entrenamiento. En el caso de clasificación binaria, esto suele formularse como la maximización del margen geométrico del clasificador \cite{van2020survey}. En aprendizaje semi-supervisado, esta idea se combina con la suposición de separación en baja densidad: además de clasificar bien los ejemplos etiquetados, se prefiere que la frontera atraviese regiones del espacio de entrada donde haya pocos datos \cite{van2020survey,pmlr-v5-goldberg09a}.

El representante más conocido es la \textit{Semi-Supervised Support Vector Machine} (S3VM), también denominada históricamente \textit{Transductive SVM} (TSVM) \cite{10.7551/mitpress/9780262033589.001.0001}. Para introducirla, repasamos primero la formulación básica de las SVM.

\subsubsection{SVM} 

Consideramos un conjunto etiquetado $\{(\mathbf{x}_i,y_i)\}_{i=1}^l$ con $\mathbf{x}_i\in\mathbb{R}^d$ e $y_i\in\{-1,1\}$, y un clasificador lineal de la forma
\begin{equation}
f(\mathbf{x})=\mathbf{w}^{\top}\mathbf{x}+b,
\end{equation}
donde $\mathbf{w}\in\mathbb{R}^d$ y $b\in\mathbb{R}$ \cite{zhu2009introduction}. Por lo que la frontera de decisión es el siguiente hiperplano:
\begin{equation}
\left\{\mathbf{x}\in\mathbb{R}^d: \mathbf{w}^{\top}\mathbf{x}+b=0\right\},
\end{equation}
mientras que la predicción se obtiene con $\mathrm{sign}(f(\mathbf{x}))$.

La distancia (no firmada) de un punto $\mathbf{x}$ a dicho hiperplano es $\frac{|f(\mathbf{x})|}{\|\mathbf{w}\|}$, y para los ejemplos etiquetados se toma la distancia con signo que es $\frac{y_i f(\mathbf{x}_i)}{\|\mathbf{w}\|}$. Con todo esto, el margen geométrico del clasificador se define como
\begin{equation}
\gamma = \min_{i=1,\dots,l} \frac{y_i\bigl(\mathbf{w}^{\top}\mathbf{x}_i+b\bigr)}{\|\mathbf{w}\|}.
\end{equation}
Primero, vamos a asumir que los datos son \textit{linealmente separables}, es decir, que existe un hiperplano tal que $y_i\bigl(\mathbf{w}^{\top}\mathbf{x}_i+b\bigr)>0$ para todo $i=1,\dots,l$. Bajo esta suposición, el \textit{hard-margin SVM} busca el hiperplano separador que maximiza $\gamma$, lo que puede escribirse como
\begin{equation}
\max_{\mathbf{w},b} \quad \min_{i=1,\dots,l} \frac{y_i\bigl(\mathbf{w}^{\top}\mathbf{x}_i+b\bigr)}{\|\mathbf{w}\|}.
\end{equation}

Este problema es invariante frente a escalado: si multiplicamos $(\mathbf{w},b)$ por una constante $c>0$, la frontera de decisión $\{\mathbf{x}: \mathbf{w}^{\top}\mathbf{x}+b=0\}$ no cambia. Por ello, podemos fijar una escala imponiendo $\min_i y_i\bigl(\mathbf{w}^{\top}\mathbf{x}_i+b\bigr)=1$, equivalente a exigir $y_i\bigl(\mathbf{w}^{\top}\mathbf{x}_i+b\bigr)\ge 1$ para todo $i$. En consecuencia, maximizar el margen es equivalente a maximizar $\frac{1}{\|\mathbf{w}\|}$, o de forma equivalente minimizar la norma de $\mathbf{w}$:
\begin{equation}
\begin{aligned}
\min_{\mathbf{w},b} \quad & \frac{1}{2}\|\mathbf{w}\|^2 \\
\text{s.a.} \quad & y_i\bigl(\mathbf{w}^{\top}\mathbf{x}_i+b\bigr) \geq 1, \quad i=1,\dots,l.
\end{aligned}
\end{equation}

Cuando los datos no son separables (esto es, no existe un hiperplano que separe perfectamente las clases), se introduce el \textit{soft margin} mediante variables de holgura $\xi_i\ge 0$ y un parámetro $C>0$ que controla el compromiso entre maximizar el margen y minimizar errores. Intuitivamente, $\xi_i$ mide cuánto viola el ejemplo $i$ la restricción de margen: $\xi_i=0$ significa que se cumple el margen, $0<\xi_i<1$ indica que el punto queda dentro del margen pero sigue estando bien clasificado, y $\xi_i>1$ corresponde a una clasificación errónea.
\begin{equation}
\begin{aligned}
\min_{\mathbf{w},b} \quad & \frac{1}{2}\|\mathbf{w}\|^2 + C \sum_{i=1}^l \xi_i \\
\text{s.a.} \quad & y_i\bigl(\mathbf{w}^{\top}\mathbf{x}_i+b\bigr) \geq 1 - \xi_i, \quad i=1,\dots,l, \\
                  & \xi_i \geq 0.
\end{aligned}
\end{equation}

Finalmente, puede expresarse en el marco de la minimización de riesgo regularizado (Sección \ref{sec:perdida_regularizacion}). Despejando $\xi_i$ de $y_i f(\mathbf{x}_i) \ge 1-\xi_i$ obtenemos $\xi_i\ge 1-y_if(\mathbf{x}_i)$, y junto con $\xi_i\ge 0$ se tiene $\xi_i=\max\bigl(0,1-y_if(\mathbf{x}_i)\bigr)$. Por lo tanto, la función objetivo se reescribe como
\begin{equation}
\min_{\mathbf{w},b} \quad \frac{1}{2}\|\mathbf{w}\|^2 + C \sum_{i=1}^l \max\bigl(0, 1 - y_i f(\mathbf{x}_i)\bigr).
\end{equation}
En este caso la función de pérdida es $c(\mathbf{x},y,f(\mathbf{x})) = \max\bigl(0, 1 - y f(\mathbf{x})\bigr)$ (pérdida \textit{hinge}) y el regularizador es $\frac{1}{2}\|\mathbf{w}\|^2$.

\subsubsection{S3VM}

La idea central de las S3VM es incorporar un conjunto no etiquetado $\{\mathbf{x}_j\}_{j=l+1}^{l+u}$ penalizando que dichos puntos queden cerca de la frontera: se prefiere que $|f(\mathbf{x}_j)|$ sea grande, de modo que la frontera se sitúe en regiones de baja densidad \cite{van2020survey,10.7551/mitpress/9780262033589.001.0001}. Una forma estándar de derivar esta penalización es introducir etiquetas latentes $\hat{y}_j\in\{-1,1\}$ para los ejemplos no etiquetados y añadir un término de pérdida \textit{hinge} $\max\bigl(0,1-\hat{y}_j f(\mathbf{x}_j)\bigr)$, que al minimizar respecto a $\hat{y}_j$ se convierte en $\max\bigl(0,1-|f(\mathbf{x}_j)|\bigr)$, conocida como \textit{hat loss} \cite{10.7551/mitpress/9780262033589.001.0001}.

Con ello, una formulación típica de S3VM es
\begin{equation}
\min_{\mathbf{w},b} \quad \frac{1}{2}\|\mathbf{w}\|^2 + C \sum_{i=1}^l \max\bigl(0, 1 - y_i f(\mathbf{x}_i)\bigr) + C^* \sum_{j=l+1}^{l+u} \max\bigl(0, 1 - |f(\mathbf{x}_j)|\bigr).
\end{equation}

En términos del marco de la Sección \ref{sec:perdida_regularizacion}, el sumatorio sobre los datos no etiquetados puede interpretarse como un término adicional dependiente de $\{\mathbf{x}_j\}$, que incorpora explícitamente la suposición de separación en baja densidad.

En la práctica suele ser necesario imponer alguna restricción adicional (por ejemplo, sobre el balance de etiquetas asignadas a los datos no etiquetados) para evitar soluciones degeneradas \cite{van2020survey}. El problema resultante es no convexo, por lo que se recurre a métodos heurísticos y aproximados; además, si la suposición de separación en baja densidad no se cumple, el uso de la información no etiquetada puede incluso degradar el rendimiento \cite{van2020survey}.

\subsection{Basados en perturbaciones}
Por otra parte, los métodos basados en perturbaciones se fundamentan en la suposición de suavidad: puntos cercanos en el espacio de entrada deberían tener etiquetas similares \cite{10.7551/mitpress/9780262033589.001.0001}. Para aprovechar esta suposición, se introducen perturbaciones o ruido estocástico en los datos de entrenamiento y se añade un término de regularización por consistencia a la función de pérdida. Este término penaliza que el modelo produzca predicciones diferentes para las versiones original y perturbada de un mismo ejemplo \cite{van2020survey}.

Este enfoque se ha popularizado especialmente en el contexto de las redes neuronales profundas, debido a su capacidad para aprender representaciones ricas y a la facilidad algorítmica de incorporar perturbaciones en el proceso de entrenamiento \cite{van2020survey}. Un ejemplo pionero es el método \textit{ladder network} \cite{rasmus2015semi}, que combina una arquitectura de red neuronal con un término de pérdida que evalúa la discrepancia entre las activaciones de la red para ejemplos limpios y perturbados.

\subsubsection{Redes neuronales}

Antes de proceder a exponer el método de las \textit{ladder networks}, repasamos brevemente el marco general de las redes neuronales, más concretamente el marco de las redes neuronales feedforward. 

Una red neuronal es un sistema que transforma un vector de entrada $\mathbf{x}\in\mathbb{R}^d$ en un vector de salida $f(\mathbf{x})\in\mathbb{R}^k$ mediante la propagación de la entrada por una serie de capas conectadas secuencialmente \cite{van2020survey}.

Cada capa de la red se compone de un conjunto de perceptrones, un modelo matemático de neurona. Cada uno de estos perceptrones recibe como entrada un vector $\mathbf{z}\in\mathbb{R}^m$, o bien el de entrada o bien las salidas de la capa anterior, y produce una salida escalar $a\in\mathbb{R}$ mediante la aplicación de una función de activación $\sigma:\mathbb{R}\to\mathbb{R}$ a una combinación lineal de las entradas, es decir, $a=\sigma(\mathbf{w}^{\top}\mathbf{z}+b)$, donde $\mathbf{w}\in\mathbb{R}^m$ y $b\in\mathbb{R}$ son los parámetros del perceptrón \cite{bishop2006pattern}.

En el contexto de aprendizaje supervisado, el entrenamiento de la red neuronal consiste en ajustar los parámetros de los perceptrones, representados por una matriz $W$ que contiene tanto los pesos $\mathbf{w}$ como los sesgos $b$, para minimizar una función de pérdida \cite{van2020survey}.

\begin{equation}
    \mathcal{L}(W) = \frac{1}{l} \sum_{i=1}^l c(\mathbf{x}_i,y_i,f(\mathbf{x}_i;W))
\end{equation}

Los parámetros se optimizan iterativamente mediante la retropropagación del error, normalmente con un método de descenso de gradiente estocástico o alguna de sus variantes \cite{NIPS2014_f033ed80}.

\subsubsection{Ladder networks}

Una vez introducidos los conceptos básicos de las redes neuronales, es posible presentar el método de las \textit{ladder networks} \cite{rasmus2015semi}. Este enfoque extiende la arquitectura de una red neuronal feedforward estándar incorporando un término de pérdida no supervisado, basado en el principio de los autoencoders eliminadores de ruido.

En concreto, el modelo se estructura en un codificador y un decodificador. Durante el entrenamiento, cada dato de entrada realiza dos pasadas por el codificador: una pasada limpia y una pasada a la que se le inyecta ruido en todas las capas ocultas. Posteriormente, el decodificador intenta ir limpiando este ruido hacia atrás, reconstruyendo las activaciones limpias originales capa por capa a partir de las representaciones perturbadas. De este modo, el término de pérdida global se define como la suma del coste de clasificación supervisado y un coste de consistencia o reconstrucción calculado a lo largo de todas las capas de la red \cite{rasmus2015semi}.

La función de pérdida total a optimizar se puede formular de la siguiente manera:
\begin{equation}
    \mathcal{L}(W) = \frac{1}{l} \sum_{i=1}^l c(\mathbf{x}_i,y_i,f(\mathbf{x}_i;W)) + \sum_{i=1}^{l+u} \sum_{k=1}^K ReconstructionLoss\bigl(\hat{\mathbf{z}}_i^k, \mathbf{z}_i^k\bigr)
\end{equation}

Donde $\mathbf{z}_i^k$ es la activación oculta de la capa $k$ en el codificador para el ejemplo limpio $\mathbf{x}_i$, mientras que $\hat{\mathbf{z}}_i^k$ es la activación reconstruida por el decodificador a partir de la versión perturbada. Finalmente, $\text{ReconstructionLoss}$ es una función que penaliza la discrepancia entre ambas (como, por ejemplo, la norma $L_2$ al cuadrado) \cite{rasmus2015semi}.

A través de esta arquitectura, las \textit{ladder networks} consiguen aprovechar simultáneamente tanto la señal de los datos etiquetados como la regularización proporcionada por los datos no etiquetados\footnote{Esta simultaneidad es, de hecho, lo que hace a las \textit{ladder networks} modelos intrínsecos, a diferencia del preentrenamiento por capas visto en la Sección \ref{sec:pre-training}, donde se usaba un \textit{encoder-decoder} como fase inicial independiente seguida de la fase supervisada.}, forzando a la red a extraer características latentes útiles y robustas al ruido que, en última instancia, mejoran la generalización del modelo predictivo \cite{rasmus2015semi}.
\section{Basados en variedades}
\label{sec: variedades}

Finalmente, los métodos basados en variedades constituyen una categoría independiente en la presente taxonomía. Por un lado, no pueden considerarse generativos, puesto que no modelan explícitamente el proceso de generación de los datos. Por otro lado, aunque (al igual que los métodos discriminativos) se centran en aprender una función de decisión, su enfoque subyacente es diferente al de las técnicas de máximo margen o de perturbaciones.

Estos métodos se apoyan en la hipótesis de la variedad: se asume que los datos observados se concentran cerca de una variedad de baja dimensión inmersa en el espacio de entrada original \cite{10.7551/mitpress/9780262033589.001.0001}. En consecuencia, el papel de los datos no etiquetados no es solo ajustar una frontera global, sino capturar la geometría de dicha variedad y orientar la búsqueda hacia funciones que respeten su estructura intrínseca \cite{van2020survey}.

\subsection{Técnicas de regularización de la variedad}

Una forma habitual de aproximar la variedad subyacente es construir un grafo de similitud entre instancias \cite{van2020survey}. Dado un conjunto de entrenamiento formado por $\{(\mathbf{x}_i, y_i)\}_{i=1}^l$ instancias etiquetadas y $\{\mathbf{x}_j\}_{j=l+1}^{l+u}$ instancias sin etiquetar, se construye un grafo no dirigido $G=(V,E)$ donde cada nodo corresponde a una instancia (etiquetada o no)\footnote{Como consecuencia, los grafos resultantes suelen tener gran tamaño cuando se dispone de una gran cantidad de instancias no etiquetadas \cite{zhu2009introduction}.}, y a cada arista $(i,j)\in E$ se le asigna un peso $w_{ij}$ que refleja la similitud entre $\mathbf{x}_i$ y $\mathbf{x}_j$. Denotamos por $W=(w_{ij})$ la matriz de pesos. La intuición es que, si $w_{ij}$ es alto, entonces es más probable que ambas instancias compartan etiqueta \cite{zhu2009introduction}.

Hay varias formas de definir la vecindad y/o la similitud. Por ejemplo, puede construirse un grafo de $k$ vecinos más cercanos (k-NN) y asignar pesos gaussianos $w_{ij} = \exp\bigl(-\frac{\|\mathbf{x}_i - \mathbf{x}_j\|^2}{2\sigma^2}\bigr)$ a las aristas \cite{zhu2009introduction}.

Una vez construido el grafo, se busca una función de decisión $f:\mathcal{X}\to\mathcal{Y}$ que no solo clasifique correctamente los ejemplos etiquetados, sino que también sea suave respecto a la estructura del grafo; es decir, que produzca salidas similares en nodos conectados por aristas con pesos altos \cite{zhu2009introduction}.

Para ello, un enfoque representativo es el presentado en \textit{On Manifold Regularization, Belkin et al., 2005} \cite{belkin2005manifold}. Partiendo de un marco de aprendizaje supervisado, se considera un kernel $K$\footnote{Un kernel es una función $K:\mathcal{X}\times\mathcal{X}\to\mathbb{R}$ que puede interpretarse como un producto escalar en un espacio de características (posiblemente de mayor dimensión) \cite{bishop2006pattern}.}, con un espacio de funciones asociado $\mathcal{H}_K$ y norma $\|\cdot\|_K$. Se busca una función $f\in\mathcal{H}_K$ que minimice la siguiente función de pérdida regularizada:
\begin{equation}
    \mathcal{L}(f) = \frac{1}{l} \sum_{i=1}^l c(\mathbf{x}_i,y_i,f(\mathbf{x}_i)) + \lambda \|f\|_K^2,
\end{equation}

donde el primer término es la pérdida de clasificación sobre los datos etiquetados y el segundo término es un regularizador que penaliza la complejidad de la función $f$ \cite{belkin2005manifold}.

Finalmente, la información no etiquetada se introduce mediante un término de regularización adicional, utilizando la matriz de pesos $W$, que penaliza la falta de suavidad de $f$ respecto a la estructura del grafo:
\begin{equation}
    \|f\|_I^2 = \frac{1}{2} \sum_{i,j} w_{ij} \bigl(f(\mathbf{x}_i) - f(\mathbf{x}_j)\bigr)^2,
\end{equation}

Este término puede interpretarse como una medida de la variación de $f$ a lo largo de las aristas del grafo, ponderada por los pesos de similitud. Además, si denotamos $n=l+u$ y definimos el vector $\mathbf{f} = \bigl(f(\mathbf{x}_1),\dots,f(\mathbf{x}_n)\bigr)^{\top}$, la matriz diagonal de grados $D$ con $D_{ii}=\sum_j w_{ij}$ y el laplaciano no normalizado $L=D-W$, entonces se obtiene una expresión matricial compacta para él mismo $\|f\|_I^2 = \mathbf{f}^{\top} L \mathbf{f}$ \cite{belkin2005manifold}.

De este modo, la función de pérdida total a minimizar es

\begin{equation}
    \mathcal{L}(f) = \frac{1}{l} \sum_{i=1}^l c(\mathbf{x}_i,y_i,f(\mathbf{x}_i)) + \lambda_A \|f\|_K^2 + \lambda_I \|f\|_I^2,
\end{equation}

donde $\lambda_A$ y $\lambda_I$ son hiperparámetros que controlan el peso relativo de cada término de regularización \cite{belkin2005manifold}.

Dependiendo de la elección del kernel $K$ y de la función de pérdida $c$, esta formulación puede dar lugar a diferentes algoritmos concretos. Un ejemplo clásico son las \textit{Laplacian SVM}, que combinan una pérdida \textit{hinge} con el término de regularización basado en el grafo. De este modo, a diferencia de las S3VM vistas en la sección anterior, la información no etiquetada se incorpora a través de la geometría del grafo y no mediante la penalización directa de puntos cercanos a la frontera de decisión \cite{van2020survey}.

En la práctica, estos métodos pueden ser costosos computacionalmente, ya que suelen requerir operar con matrices de tamaño $n\times n$ (donde $n=l+u$). En particular, en algunos casos el coste puede alcanzar $O(n^3)$, o bien $O(c\,n^2)$ para un $c \ll n$ \cite{van2020survey}. Además, el rendimiento es muy sensible a cómo se construye el grafo y a la elección y escala de los pesos $w_{ij}$; si el grafo no refleja bien la geometría de los datos, la información no etiquetada puede propagarse de forma perjudicial \cite{zhu2009introduction}.